% ---
% title: A constant density on a triangle
% tags: joint-continuous, marginal, independent-rv
% ---

\begin{problem}
	The random point $(X,Y)$ is uniformly distributed on the triangle
	\[ T=\{(x,y): 0<y<x<1\}: \]
	the joint density is a constant $c$ on $T$ and $0$ elsewhere,
	\[
		f(x,y)=\begin{cases}
			c & 0<y<x<1,   \\
			0 & \text{otherwise.}
		\end{cases}
	\]
	(a) Find $c$.
	(b) Find the marginal densities $f_X$ and $f_Y$.
	(c) Find $P(X<\frac12)$ and $P(Y>\frac12)$.
	(d) Are $X$ and $Y$ independent?
\end{problem}

\begin{solution}
	\idea A joint density integrates to $1$ over the plane. When the density is
	a constant $c$ on a region, the integral is $c$ times the area of the
	region, so probabilities are $c\,\times$ area. A marginal density is found
	by integrating out the other variable; for this we must read off from the
	picture where the other variable can be, for a fixed value of the first.

	\step{The region.} The condition $0<y<x<1$ means: $x$ is between $0$ and
	$1$, and $y$ is between $0$ and $x$. This is the triangle below the
	diagonal $y=x$ of the unit square, with corners $(0,0)$, $(1,0)$, $(1,1)$.
	\begin{center}
		\begin{tikzpicture}[scale=4,>=stealth]
			\fill[teal!15] (0,0)--(1,0)--(1,1)--cycle;
			\draw[navy,thick] (0,0)--(1,0)--(1,1)--cycle;
			\draw[->] (-0.05,0)--(1.15,0) node[right] {$x$};
			\draw[->] (0,-0.05)--(0,1.15) node[above] {$y$};
			\draw[gray,dashed] (0,1)--(1,1);
			\node[below] at (1,0) {$1$};
			\node[left] at (0,1) {$1$};
			\node[below left] at (0,0) {$0$};
			% vertical slice at x = 0.4
			\draw[navy,very thick] (0.4,0)--(0.4,0.4);
			\node[below] at (0.4,0) {$x$};
			\node[navy,right,font=\footnotesize] at (0.4,0.13) {$0<y<x$};
			% horizontal slice at y = 0.75
			\draw[gold!80!black,very thick] (0.75,0.75)--(1,0.75);
			\draw[gray,dotted] (0,0.75)--(0.75,0.75);
			\node[left] at (0,0.75) {$y$};
			\node[gold!60!black,right,font=\footnotesize] at (1,0.75) {$y<x<1$};
			\node[rotate=45,above,font=\footnotesize] at (0.55,0.55) {$y=x$};
			% the point (0.2, 0.5)
			\fill[red!70!black] (0.2,0.5) circle (0.4pt);
			\node[red!70!black,above,font=\footnotesize] at (0.2,0.5) {$(0.2,\,0.5)$};
			\node[navy] at (0.85,0.4) {$T$};
		\end{tikzpicture}
	\end{center}
	The navy segment shows where $y$ can be when $x$ is fixed; the gold segment
	shows where $x$ can be when $y$ is fixed. (The red point is used in (d).)

	\step{(a)} The total probability must be $1$:
	\[ \iint f(x,y)\,dx\,dy=c\cdot\text{area}(T)=1. \]
	The triangle has base $1$ and height $1$, so its area is
	$\frac12\cdot1\cdot1=\frac12$. Hence $c\cdot\frac12=1$ and
	\[ c=2. \]
	The same with a double integral: for fixed $x$ the inner integral runs over
	$0<y<x$,
	\[
		\int_0^1\!\left(\int_0^x c\,dy\right)dx=\int_0^1 cx\,dx
		=c\left[\frac{x^2}{2}\right]_0^1=\frac c2 .
	\]

	\step{(b)} \emph{Marginal of $X$.} Fix $x$ with $0<x<1$ and integrate over
	$y$. The density is $2$ only for $0<y<x$ (the navy segment), and $0$ for
	other $y$:
	\[
		f_X(x)=\int_{-\infty}^{\infty}f(x,y)\,dy=\int_0^x 2\,dy
		=\bigl[2y\bigr]_0^x=2x-0=2x,\qquad 0<x<1,
	\]
	and $f_X(x)=0$ outside $(0,1)$.

	\emph{Marginal of $Y$.} Fix $y$ with $0<y<1$ and integrate over $x$. The
	density is $2$ only for $y<x<1$ (the gold segment):
	\[
		f_Y(y)=\int_{-\infty}^{\infty}f(x,y)\,dx=\int_y^1 2\,dx
		=\bigl[2x\bigr]_y^1=2-2y=2(1-y),\qquad 0<y<1,
	\]
	and $f_Y(y)=0$ outside $(0,1)$.

	Check: $\int_0^1 2x\,dx=\bigl[x^2\bigr]_0^1=1$ and
	$\int_0^1 2(1-y)\,dy=\bigl[2y-y^2\bigr]_0^1=2-1=1$. Both are densities.
	This makes sense: $X$ is the larger coordinate, so it tends to be large
	($f_X$ grows); $Y$ is the smaller one, so it tends to be small ($f_Y$
	falls).

	\step{(c)} With the marginal density of $X$:
	\[
		P\Bigl(X<\tfrac12\Bigr)=\int_0^{1/2}2x\,dx=\bigl[x^2\bigr]_0^{1/2}
		=\frac14-0=\frac14 .
	\]
	With the marginal density of $Y$ (an antiderivative of $2(1-y)$ is
	$-(1-y)^2$):
	\[
		P\Bigl(Y>\tfrac12\Bigr)=\int_{1/2}^{1}2(1-y)\,dy
		=\Bigl[-(1-y)^2\Bigr]_{1/2}^{1}=0-\Bigl(-\frac14\Bigr)=\frac14 .
	\]
	Geometry gives the same faster. The part of $T$ with $x<\frac12$ is the
	triangle with corners $(0,0)$, $(\frac12,0)$, $(\frac12,\frac12)$; the part
	with $y>\frac12$ is the triangle with corners $(\frac12,\frac12)$,
	$(1,\frac12)$, $(1,1)$. Both have legs $\frac12$, so area
	$\frac12\cdot\frac12\cdot\frac12=\frac18$, and probability
	$2\cdot\frac18=\frac14$.

	\step{(d)} $X$ and $Y$ are independent exactly when
	$f(x,y)=f_X(x)f_Y(y)$ for all $x,y$ (apart from a set of zero area). Take
	the point $(0.2,\,0.5)$. Here $y>x$, so it lies outside the triangle and
	\[ f(0.2,\,0.5)=0. \]
	But both marginals are positive there:
	\[ f_X(0.2)\,f_Y(0.5)=(2\cdot0.2)\cdot2(1-0.5)=0.4\cdot1=0.4>0. \]
	The product rule fails, and it fails on the whole region above the
	diagonal (inside the unit square), which has area $\frac12$, not just at one
	point. So $X$ and $Y$ are \emph{not} independent.

	In words: knowing $X$ changes what $Y$ can be. Given $X=0.2$, the value of
	$Y$ must be below $0.2$, although on its own $Y$ can be anywhere in
	$(0,1)$.

	\begin{answer}
		(a) $c=2$ \quad
		(b) $f_X(x)=2x$ for $0<x<1$, $f_Y(y)=2(1-y)$ for $0<y<1$ (and $0$
		elsewhere) \quad
		(c) $P(X<\frac12)=\frac14$, $P(Y>\frac12)=\frac14$ \quad
		(d) No, they are not independent.
	\end{answer}
	\begin{remark}
		\textbf{Common mistake:} writing the bounds $0<y<1$ and $0<x<1$ in the
		marginal integrals, as for the unit square. The bounds of the inner
		integral come from the region and may depend on the variable held
		fixed: here $y$ runs up to $x$, and $x$ starts at $y$.

		A quick test: if the region where $f>0$ is not a rectangle (with sides
		parallel to the axes), the random variables cannot be independent. The
		product $f_X(x)f_Y(y)$ is positive on a whole rectangle, the set of all
		$(x,y)$ with $f_X(x)>0$ and $f_Y(y)>0$, so the joint density would have
		to be positive there too.
	\end{remark}
\end{solution}
